\documentclass[11pt]{article}
\usepackage[a4paper,margin=1in]{geometry}
\usepackage{fontspec}
\usepackage[boldfont=false]{xeCJK}
\setCJKmainfont{FandolSong-Regular.otf}
\usepackage{amsmath}
\usepackage{amssymb}
\usepackage{array}
\usepackage{booktabs}
\usepackage{graphicx}
\usepackage{adjustbox}
\usepackage{enumitem}
\setlist[itemize]{topsep=2pt,partopsep=0pt,parsep=0pt,itemsep=2pt,leftmargin=1.4em}
\setlist[enumerate]{topsep=2pt,partopsep=0pt,parsep=0pt,itemsep=2pt,leftmargin=1.6em}
\usepackage{parskip}
\usepackage{fvextra}
\fvset{breaklines=true,breakanywhere=true,fontsize=\small}
\setlength{\parindent}{0pt}

\begin{document}

\begin{center}
  {\LARGE\bfseries Human influences on ecosystems}\\[0.35em]
  {\large Igcse Biology}
\end{center}
\vspace{0.6em}

\section*{Increasing food supply}

To grow more food, humans use:

\begin{itemize}
\item \textbf{agricultural machinery} to farm larger areas of land quickly.

\item chemical \textbf{fertilisers} 化肥 to improve \textbf{yields} 产量.

\item \textbf{insecticides} 杀虫剂 to kill insect pests.

\item \textbf{herbicides} 除草剂 to kill \textbf{weeds} 杂草 that compete with the crops.

\item \textbf{selective breeding} 选择育种 to improve \textbf{crop plants} 农作物 and \textbf{livestock} 牲畜.

\end{itemize}

A \textbf{monoculture} 单一栽培 (growing one crop over a large area) gives a big, easy harvest, but one pest or disease can destroy the whole crop, and it lowers \textbf{biodiversity} 生物多样性. Intensive livestock farming (many animals in a small space) gives cheap meat, but raises worries about animal welfare and the spread of disease.

\section*{Habitat destruction}

Biodiversity is the number of different \textbf{species} 物种 living in an area. Humans destroy \textbf{habitats} 栖息地 (the places where organisms live) by:

\begin{itemize}
\item clearing land for housing, crops and livestock.

\item extracting natural \textbf{resources} 资源, such as by mining.

\item \textbf{pollution} 污染 of fresh water and the sea.

\end{itemize}

By changing \textbf{food webs} 食物网 and \textbf{food chains} 食物链, humans harm habitats.

\textbf{Deforestation} 森林砍伐 (cutting down forests) is a major example. It causes:

\begin{itemize}
\item lower biodiversity and the \textbf{extinction} 灭绝 of species.

\item loss of soil (no roots to hold it, so it washes away).

\item more flooding.

\item more \textbf{carbon dioxide} 二氧化碳 in the air, because there are fewer trees to take it in.

\end{itemize}

\section*{Pollution}

\subsection*{Water pollution}

Untreated \textbf{sewage} 污水 and excess fertiliser washed into rivers and lakes damage water \textbf{ecosystems} 生态系统.

\textbf{(Supplement)} Excess fertiliser causes \textbf{eutrophication} 富营养化:

\begin{enumerate}
\item the fertiliser adds \textbf{nitrate ions} 硝酸根离子 and other ions to the water.

\item \textbf{producers} 生产者 such as algae grow very fast.

\item they soon die, and \textbf{decomposition} 分解 increases.

\item \textbf{decomposers} 分解者 use up the \textbf{oxygen} 氧气 in \textbf{aerobic respiration} 有氧呼吸.

\item the dissolved oxygen falls, so fish and other organisms die.

\end{enumerate}

\subsection*{Plastics and air pollution}

\textbf{Non-biodegradable} 不可降解 \textbf{plastics} 塑料 do not rot away. They build up in the sea and on land, harming wildlife — animals may eat them or get trapped.

\textbf{Methane} 甲烷 and carbon dioxide are greenhouse gases. Too much of them strengthens the \textbf{greenhouse effect} 温室效应 (the enhanced greenhouse effect), which causes \textbf{climate change} 气候变化.

\section*{Conservation}

\subsection*{Sustainable resources}

A \textbf{sustainable} 可持续 resource is produced as fast as it is used up, so it never runs out. Forests and \textbf{fish} 鱼类 stocks can both be managed in this way.

\textbf{(Supplement)} Forests are \textbf{conserved} 保护 by education, protected areas, \textbf{quotas} 配额 (limits on how much is taken) and replanting. Fish stocks are conserved by \textbf{closed seasons} 禁渔期, protected areas, controlled net mesh sizes (so young fish can escape), quotas and monitoring.

\subsection*{Saving endangered species}

A species becomes \textbf{endangered} 濒危 or extinct when its numbers fall, because of climate change, habitat destruction, \textbf{hunting} 狩猎, \textbf{overharvesting} 过度捕捞, pollution and \textbf{introduced species} 外来物种.

Endangered species can be saved by:

\begin{itemize}
\item monitoring and protecting the species and its habitat.

\item education.

\item \textbf{captive breeding} 圈养繁殖 programmes, which breed animals in zoos. \textbf{(Supplement)} These may use \textbf{artificial insemination} 人工授精 (AI) or \textbf{in vitro fertilisation} 体外受精 (IVF).

\item \textbf{seed banks} 种子库, which store seeds safely.

\end{itemize}

\textbf{(Supplement)} We conserve nature to keep up biodiversity, reduce extinction, protect ecosystems, and keep ecosystem functions going — such as recycling nutrients and providing food, \textbf{drugs} 药物, fuel and useful \textbf{genes} 基因. If a population becomes too small, it loses \textbf{genetic variation} 变异, which makes it harder for the species to survive future change.

\section*{Exam tips}

\begin{itemize}
\item More food: machinery, fertilisers, insecticides, herbicides, selective breeding. Monocultures and intensive farming each have advantages and disadvantages.

\item Deforestation: less biodiversity, extinction, soil loss, flooding, more carbon dioxide.

\item Eutrophication (Supplement): fertiliser → more producers → more decomposers → less dissolved oxygen → organisms die.

\item Greenhouse gases (carbon dioxide, methane) → enhanced greenhouse effect → climate change.

\item Conservation: sustainable use, protected areas, quotas, captive breeding and seed banks. Small populations lose genetic variation.

\end{itemize}


\end{document}
